%============================================================
%  Chapter 4 – System Architecture and Design
%============================================================
\chapter{System Architecture and Design}

\section{Architectural Overview}

\projectname{} follows a \textbf{three-tier client-server architecture} with a clear separation of concerns between the presentation, application logic, and data tiers. The system is designed as a RESTful API backend consumed by a React single-page application frontend.

\begin{figure}[H]
\centering
\setlength{\fboxsep}{10pt}
\fbox{%
\begin{minipage}{0.85\textwidth}
\begin{center}
\textbf{Three-Tier Architecture}\\[0.5cm]
\begin{tabular}{|c|}
\hline
\\
\textbf{Tier 1 – Presentation Layer} \\
React + Vite SPA (TypeScript) \\
React Router, TanStack Query, Shadcn UI \\
{\small Deployed on user's browser} \\[0.2cm]
\hline
\\
\textbf{Tier 2 – Application Layer} \\
Node.js + Express.js REST API (TypeScript) \\
Controllers -- Services -- Utilities \\
Redis Cache, Nodemailer, Cloudinary, OpenAI \\[0.2cm]
\hline
\\
\textbf{Tier 3 – Data Layer} \\
MongoDB (document store) + Prisma ORM \\
Redis (key-value cache and session data) \\[0.2cm]
\hline
\end{tabular}
\end{center}
\end{minipage}%
}
\caption{High-Level Three-Tier Architecture of \projectname{}}
\label{fig:architecture}
\end{figure}

\section{Technology Stack}

\subsection{Backend}

\begin{table}[H]
\centering
\caption{Backend Technology Stack}
\begin{tabularx}{\textwidth}{llX}
\toprule
\textbf{Technology} & \textbf{Version} & \textbf{Purpose} \\
\midrule
Node.js    & 20.x  & JavaScript/TypeScript server runtime. \\
Express.js & 5.x   & HTTP web framework providing routing, middleware, and request handling. \\
TypeScript & 5.x   & Typed superset of JavaScript providing compile-time safety. \\
Prisma ORM & 6.x   & Type-safe ORM for database access and schema management. \\
MongoDB    & 7.x   & NoSQL document database for flexible, scalable data storage. \\
Redis      & 7.x   & In-memory data store for caching, rate-limiting, and OTP storage. \\
Cloudinary & 2.x   & Cloud-based image and file storage with CDN delivery. \\
Nodemailer & 8.x   & SMTP email sending for OTP, password reset, and notifications. \\
OpenAI     & 6.x   & GPT API integration for the AI project scoping assistant. \\
Passport.js & 0.7  & Authentication middleware for Google OAuth 2.0. \\
bcrypt.js  & 3.x   & Password hashing library. \\
jsonwebtoken & 9.x & JWT signing and verification. \\
Zod        & 4.x   & Runtime schema validation for request bodies. \\
Multer     & 2.x   & Multipart form-data middleware for file uploads. \\
node-cron  & 4.x   & Monthly SkillToken award cron job scheduling. \\
\bottomrule
\end{tabularx}
\label{tab:backend_stack}
\end{table}

\subsection{Frontend}

\begin{table}[H]
\centering
\caption{Frontend Technology Stack}
\begin{tabularx}{\textwidth}{llX}
\toprule
\textbf{Technology} & \textbf{Version} & \textbf{Purpose} \\
\midrule
React       & 18.x & Component-based UI framework. \\
Vite        & 5.x  & Next-generation frontend build tool and dev server. \\
TypeScript  & 5.x  & Static typing for the frontend codebase. \\
React Router & 6.x & Client-side routing and navigation. \\
TanStack Query & 5.x & Server state management, caching, and synchronisation. \\
Axios       & 1.x  & HTTP client for API communication. \\
Shadcn/Radix UI & -- & Accessible, composable UI component primitives. \\
Tailwind CSS & 3.x & Utility-first CSS framework for styling. \\
Lucide React & -- & Icon library. \\
React Hook Form & -- & Performant, accessible form management. \\
\bottomrule
\end{tabularx}
\label{tab:frontend_stack}
\end{table}

\section{Backend Application Structure}

The backend follows a layered architecture within a monorepo structure. The \texttt{src/} directory is organised as follows:

\begin{lstlisting}[caption={Backend Directory Structure}]
src/
+-- app.ts               # Express app: middleware + route registration
+-- server.ts            # HTTP server bootstrap + port binding
+-- config/              # Configuration modules
|   +-- cloudinary.ts    # Cloudinary SDK init
|   +-- passport.ts      # Google OAuth strategy definition
|   +-- prisma.ts        # Prisma client singleton
|   \-- redis.ts         # Redis client configuration
+-- controllers/         # Route handler functions (thin layer)
|   +-- auth.controller.ts
|   +-- freelancer-onboarding.controller.ts
|   +-- freelancer.controller.ts
|   +-- project.controller.ts
|   +-- proposal.controller.ts
|   +-- token.controller.ts
|   +-- tracking.controller.ts
|   +-- contract.controller.ts
|   \-- metadata.controller.ts
+-- routes/              # Express Router definitions
|   +-- auth.routes.ts
|   +-- freelancer.routes.ts
|   +-- project.routes.ts
|   +-- proposal.routes.ts
|   +-- token.routes.ts
|   +-- tracking.routes.ts
|   +-- category.routes.ts
|   +-- metadata.routes.ts
|   \-- ai.routes.ts
+-- middlewares/         # Express middleware functions
|   +-- auth.middleware.ts   # JWT verification
|   +-- role.ts              # Role-based access guard
|   +-- upload.middleware.ts # Multer file handling
|   \-- logger.middleware.ts # Request logging
+-- services/
|   +-- token.service.ts
|   \-- tracking.service.ts
+-- modules/
|   \-- browse/          # Browse and discovery module
+-- utils/
|   +-- email.ts             # Email templates
|   +-- jwt.ts               # Token helpers
|   +-- redis.ts             # Redis wrappers
|   +-- uploadToCloudinary.ts
|   +-- profileCompletion.ts
|   +-- tokenCalculator.ts
|   \-- validators.ts
\-- ai/                  # AI assistant module
\end{lstlisting}

\section{Request Lifecycle}

Every HTTP request to the \projectname{} API follows a deterministic lifecycle:

\begin{enumerate}
  \item \textbf{CORS Check}: The CORS middleware validates the \texttt{Origin} header against the allowed frontend list.
  \item \textbf{Body Parsing}: \texttt{express.json()} and \texttt{cookie-parser} decode the request payload and cookies.
  \item \textbf{Request Logging}: The custom \texttt{requestLogger} middleware records method, path, and response time.
  \item \textbf{Authentication}: The \texttt{protect} middleware verifies the JWT access token from the request header and attaches the decoded user payload to \texttt{req.user}.
  \item \textbf{Authorisation}: Role-specific middleware (e.g., \texttt{requireRole('CLIENT')}) checks that the authenticated user has the required role.
  \item \textbf{Controller}: The matching route handler (controller function) is called. Zod schema validation may occur here.
  \item \textbf{Service/ORM}: The controller calls Prisma to query or mutate the database, or delegates to a service function.
  \item \textbf{Response}: The controller sends a JSON response with the appropriate HTTP status code.
\end{enumerate}

\section{Frontend Application Structure}

The React frontend is organised by feature and role:

\begin{lstlisting}[caption={Frontend Directory Structure}]
src/
+-- App.tsx              # Root component: providers + routing tree
+-- main.tsx             # Vite entry point
+-- contexts/
|   +-- AuthContext.tsx  # Global auth state (user, token, logout)
|   \-- ThemeContext.tsx # Light/dark theme toggle
+-- components/
|   +-- ProtectedRoute.tsx   # Auth + role guard components
|   +-- layout/              # Navbar, sidebar, footer
|   +-- common/              # Shared components (cards, badges)
|   +-- browse/              # Browse-specific components
|   +-- client/              # Client-specific components
|   +-- modals/              # Modal dialogs
|   \-- ui/                  # Shadcn UI primitives
+-- pages/
|   +-- LandingPage.tsx
|   +-- AboutPage.tsx
|   +-- HowItWorksPage.tsx
|   +-- auth/               # Login, Signup, OTP, OAuth, Reset
|   +-- client/             # 14 client-side pages
|   +-- freelancer/         # 12 freelancer-side pages
|   +-- admin/              # 6 admin pages
|   \-- legal/              # Privacy policy, Terms of service
+-- hooks/                  # Custom React hooks
\-- lib/                    # Utility functions, API client config
\end{lstlisting}

\section{Routing Architecture}

The frontend uses React Router v6 with a nested route structure implementing three route guard types:

\begin{itemize}
  \item \textbf{Public Routes}: Accessible without authentication (Landing, About, Login, Signup, etc.).
  \item \textbf{Guest Routes}: Only accessible when \textit{not} authenticated; redirect logged-in users to their dashboard.
  \item \textbf{Protected Routes}: Require authentication. Further subdivided by \texttt{RoleRoute} into CLIENT-only, FREELANCER-only, and ADMIN-only sub-trees.
\end{itemize}

\section{State Management}

\begin{table}[H]
\centering
\caption{Frontend State Management Strategy}
\begin{tabularx}{\textwidth}{lX}
\toprule
\textbf{State Type} & \textbf{Mechanism} \\
\midrule
Authentication state  & React Context (\texttt{AuthContext}) — persisted in memory + HTTP-only cookie for refresh token. \\
Server state (API data) & TanStack Query — with automatic caching, refetching, and invalidation. \\
Form state            & React Hook Form — with Zod schema validation. \\
Theme preference      & React Context (\texttt{ThemeContext}) — persisted in \texttt{localStorage}. \\
UI state (modals, etc.) & Local \texttt{useState} within components. \\
\bottomrule
\end{tabularx}
\end{table}
